\section{Comparison protocols as bounded-degree relations}
\label{sec:protocol}

We give the comparison implementation used in the threshold gadgets.
It is the half-splitting comparison underlying Nisan's
protocol~\cite{Nisan93}; explicit parameters are useful for applying
\HTC. Fix the circuit scale $S$, and suppose $t\le S$ parties hold
nonnegative integers $X_i$, each of bit length polynomial in $S$.
The public target $T$ has the same bit bound. We compare
$\sum_iX_i>T$. Targets outside the possible sum range are handled by
public constant answers. Pad the bit width to a power of two $B'$ and
put $J=\log_2 B'=O_a(\log S)$ when the input description is at most
$S^a$ bits.

\subsection{The half-splitting step}

For a current even width $B$, put $d=B/2$ and write
\[
 X_i=2^dH_i+L_i,\quad T=2^dT_H+T_L,
 \quad 0\le L_i,T_L<2^d.
\]
If $\sum_iH_i=T_H-h$ for some $h\in\{0,\ldots,t-1\}$, then the
comparison is equivalent to
\begin{equation}\label{eq:low-comparison}
 \sum_i L_i>T_L+h2^d.
\end{equation}
If there is no such $h$, put $u=\sum_iH_i-T_H$. Then either $u\ge1$
or $u\le-t$. In the former case the original comparison is true, since
$2^du+\sum_iL_i-T_L\ge1$; in the latter case it is false. Both
answers equal the answer to $\sum_iH_i>T_H$. Thus one may recurse on
that high-half comparison. The range for $h$ includes $t-1$.

Let $\delta=2^{-h_0}$, with integer $h_0\ge1$. Choose the least power
of two $K$ such that
\begin{equation}\label{eq:prime-family}
 K\ge 4JS(B'+\lceil\log_2S\rceil+3)\,2^{h_0}.
\end{equation}
Use the $K$ smallest primes greater than $S$. At each level a fresh
$\log_2K$-bit seed chooses one such prime $p$. The parties broadcast
$H_i\bmod p$, and public arithmetic computes the least nonnegative
residue
\[
 h'=T_H-\sum_i(H_i\bmod p)\pmod p.
\]
If $h'<t$, follow~\eqref{eq:low-comparison} with $h=h'$; otherwise
follow the high-half comparison. Since $p>t$, a genuine $h$ is always
returned correctly. When no genuine $h$ exists, an incorrect choice
requires $p$ to divide one of the nonzero integers
\[
 \sum_iH_i-T_H+j,\qquad 0\le j<t.
\]
After the range checks, each has bit length at most
$B'+\lceil\log_2S\rceil+3$. An integer of that bit length has at most
that many distinct prime factors. Consequently the conditional error
at a level is at most $\delta/(4J)$, even after fixing all earlier
messages and seeds. A union bound on the earliest erroneous level gives
error at most $\delta$. At width one the parties send their remaining
bits and compare exactly.

The prime bound of Rosser and Schoenfeld~\cite{Rosser62} implies that
the largest prime used has
\begin{equation}\label{eq:prime-bits}
 \pi=O_a(h_0+\log S)
\end{equation}
bits. Each seed is only a prime index. For a particular index,
enumerate integers and use trial division until reaching that prime.
The polynomial upper bound on the indexed prime implies
$2^{O_a(h_0+\log S)}$ time and polynomial bit space. This construction
stores the selected primes, rather than the entire family. The total
seed length is
\begin{equation}\label{eq:seedlength}
 R=J\log_2K=O_a(h_0\log S+\log^2S).
\end{equation}
Fixed-width message frames, range checks, and forced zero padding
after an early stop make every transcript canonical.

\subsection{Two tensor realizations}

Fix the seed. If a party reads at most $b$ original input bits and
the full transcript has $C$ bits, make one factor for each party,
reading its private input and a copy of the transcript. The factor
checks only that party's messages, each against the earlier prefix.
A referee factor checks public control, formatting, padding, and the
output. Sharing transcript bits by equality trees gives
\begin{equation}\label{eq:full-transcript}
 M=O(tb+tC+C+1),\qquad
 \Delta\le\max\{b+C,C+1,3\}.
\end{equation}
For fixed private inputs, the initial message is determined. Inductively
each next message, stop flag, and final output is determined; the
checks force exactly this execution. Hence the gadget is a functional
relation even for a seed at which the comparison is wrong.

For higher precision we use the addition-specific compact
realization. Encode each public branch (a value of $h$ or ``no such
value'') in $\lambda=\lceil\log_2(t+1)\rceil$ bits. The factor for
party $i$ reads its original input, the public branch history, its own
$J$ residues, and its final remaining bit. It replays its local
state and checks these messages. It reads no other party's residues.
At each level a public factor uses earlier history to produce the
target residue and the stop flag. A fan-in-two Boolean circuit adds
the $t$ residues modulo the known prime, subtracts from the target
modulo that prime, and compares the result with $t$. Addition of two
$\pi$-bit residues, comparison with $p$, subtraction, and multiplexing
have $O(\pi)$ gates and wires. Therefore this public arithmetic costs
$O(t\pi)$ wires per level. A final factor checks the output from the
$t$ last bits and the public history.

Counting original input ports, residues, arithmetic wires, and history
copies gives
\begin{align}
 M&=O\bigl(tb+tJ\pi+tJ^2\lambda+J\pi+1\bigr),
       \label{eq:compact-edges}\\
 \Delta&=O\bigl(\max\{b+J(\pi+\lambda),\,
 \pi+J\lambda,\,t+J\lambda,\,3\}\bigr).
       \label{eq:compact-degree}
\end{align}
The history-copy bound is a convenient overestimate of the actual
port count. Given the inputs, earlier branches determine current
messages; those messages determine the next branch. This induction
also proves uniqueness for the compact realization. All local entry
programs use polynomial-time integer arithmetic on the original
weights and the selected primes.

\subsection{Shared seeds and the inverse-polynomial CAPP bound}

For a circuit, suppose gate $v$ is replaced by $f_{v,\rho}$ satisfying
\[
 \Pr_\rho[f_{v,\rho}(y)\ne f_v(y)]\le\eta_v
 \quad\hbox{for each fixed local input }y.
\]
Every gate may read a prefix of the same uniform seed. For an original
input $x$, evaluate the original circuit and let $y_v(x)$ be its fixed
local input at gate $v$. If the replaced circuit disagrees, its earliest
erroneous gate has input exactly $y_v(x)$. Therefore
\begin{equation}\label{eq:gate-error}
 \Pr_\rho[C_\rho(x)\ne C(x)]
 \le\sum_v\Pr_\rho[f_{v,\rho}(y_v(x))\ne f_v(y_v(x))]
 \le\sum_v\eta_v.
\end{equation}
No independence between gates is needed.

For clarity, the inverse-polynomial case already follows from the
full-transcript realization. Put $b=\lceil S^{2/3}\rceil$ and leave
smaller gates exact. For a large threshold gate, replace negative
weights by literals: with $u_j=|w_j|$ and
$\xi_j=x_j$ for $w_j\ge0$, $\xi_j=1-x_j$ otherwise,
\[
 \sum_jw_jx_j\ge\theta
 \quad\Longleftrightarrow\quad
 \sum_j u_j\xi_j>\theta+\sum_{w_j<0}|w_j|-1.
\]
Party $i$ computes the nonnegative sum in one block of at most $b$
inputs. Give each large gate error $\eta=\eps/(2S)$, using the
protocol with a dyadic error at most this value. If $\eps\ge S^{-a}$,
then $R=O_a(\log^2S)$ and a gate with fan-in $r_v$ has
$t_v\le2r_v/b$ parties and transcript
$C_v=O_a(t_v\log^2S)$. The total additional edges in
\eqref{eq:full-transcript} satisfy
\[
 \sum_{r_v>b}O(t_vC_v)
 \le O_a\left(\frac{\log^2S}{b^2}\sum_vr_v^2\right)
 \le O_a(S^{2/3}\log^2S)=o(S).
\]
The degree is at most
$S^{2/3}+O_a(S^{1/3}\log^2S)+O(1)=o(S^{9/10})$.
For each fixed seed, the circuit has a unique-extension network and
hence an exact count $Z_\rho$ from \HTC\ with $P_0=n+1$.
Enumerating all seeds gives the rational
\begin{equation}\label{eq:capp-output}
 \widehat p=\frac{\sum_\rho Z_\rho}{2^{n+R}}.
\end{equation}
There are at most $S$ gates, so~\eqref{eq:gate-error} bounds the
error by $\eps$. The seed enumeration is absorbed by
$2^{O_a(g(S))}$, and the accumulator has $n+R+1$ bits. This proves
the inverse-polynomial assertion of Theorem~\ref{thm:capp} under
\HTC. The compact realization supplies its stated endpoint in the
next section.
